\section{関連研究}
\subsection{創作活動におけるAIの活用}
AI技術の急速な発展に伴い, AIを創作活動に応用する研究が盛んに行われている. その分野は多岐に渡る. 

Ko ら\cite{ref:Ko23}は, ビジュアルアーティストがLTGM(Large-scale Text-to-image Generation Model)の導入についてどのような理解を持っているのかを調査し, LTGMを活用するユーザーインタフェースの設計に有用なデザインガイドラインを提示した. Louie ら\cite{ref:Louie20}は, AI の楽曲生成が非決定的に行われる問題に対処するため, 生成される楽曲を特定の声に制限し, 生成物の類似性を制御するセマンティックスライダーなどを実装したツールを提案した. Guzdial ら\cite{ref:Guzdial19}は, ゲームレベルデザイナーのパートナーとなるようなAIエージェント, Morai Makerを開発した.

このように, AI の創造的活動への適応は, 画像や音楽, ゲームデザインにまで広がっている.
\subsection{言語モデルとストーリー創作支援}
%言語モデルは, Trasformer\cite{ref:Vaswani17}の登場を起点とし, 著しい発展を遂げてきた. Transformer をベースとしたモデルには BERT\cite{ref:Devlin18} や GPT\cite{ref:Radford20} などがある. 学習データとモデルの大規模化により高い汎化性能を獲得し, GPTは, その系譜を GPT-4\cite{ref:OpenAI23}にまで進歩させている. 
言語モデルの急速な発展に伴い, これらをストーリークリエーターの創作支援に活用するシステムの研究が進められている.

大曽根ら\cite{ref:Osone21}は日本の小説家に用いられることを想定したAIシステムBunchoを開発した. Bunchoは, 日本語記事の大規模データセットと web 上の小説によって学習された GPT-2\cite{ref:Radford20} を利用しており, キーワードから, タイトルや比較的短いあらすじを生成することができる. Yuanら\cite{ref:Yuan22}は, LaMDA\cite{ref:Thoppilan22} と物語の共創が行えるインターフェース, Wordcraft を開発した. Wordcraft では, 自身で書いた文章の続きを書かせることや, 文章の一部を指定して, その部分の代替案を出力させることができる. これらのシステムは, 有意義な共創体験を与えるという評価を得た一方で, ストーリーに対する認識の欠如や生成物の一貫性が不足しているという評価を受けた.

%利用した作家からは, AIとの共創体験が楽しいものであったという評価を受けた. 一方で, 生成物の正確さや一貫性という観点では, 人の手で書かれたものよりも有意に低い評価を受けた.

 %AI との対話によって文章を改変していくことができるシステムであるが, AI の返答がうまく文脈を認識していないと感じられる課題が生じている. 

日笠ら\cite{ref:Higasa23}は, 物語構造を用いGPT-3\cite{ref:Brown20} を制御することで生成したプロットユニットを, 意味的類似度で繋ぎ合わせプロットの自動生成を行なった. ここでいうプロットとは, 物語の大きな流れを示す文章のことであり, プロットユニットとは, プロットを展開ごとに分割した要素のことを表す. この研究により, プロット生成における物語構造の有用性が示唆された. しかしながら, このシステムはクリエーターが直接操作できず, クリーエーターの意思を反映することができない.

%日笠ら\cite{ref:Higasa23}は, 物語構造を用いて GPT-3\cite{ref:Brown20} を制御し, プロットユニットを生成をすることで, プロットの自動生成に取り組んだ. ここでいうプロットとは, 物語の大きな流れを示す文章のことであり, プロットユニットとは, プロットを展開ごとに分割した要素のことを指す. 個別に生成したプロットユニットを意味的類似度で繋ぎ合わせ, 一つの大きなプロットを作り上げた. このことにより, プロット生成において, 物語構造の活用が有用であることを示唆した. しかしながら, このシステムはストーリークリエーターが自身の手で動かすことのできるシステムではなく, クリーエーターの意思を反映することができない.